\documentclass[11pt,a4paper]{article}

% ---------------------------------------------------------------------------
% Preliminary Notes on Infrastructure, Data Standards and Data Management
% for Federated Learning in Biological Countermeasures
%
% Build:
%   pdflatex fl-biocountermeasures-preliminary
%   bibtex   fl-biocountermeasures-preliminary
%   pdflatex fl-biocountermeasures-preliminary
%   pdflatex fl-biocountermeasures-preliminary
% ---------------------------------------------------------------------------

\usepackage[utf8]{inputenc}
\usepackage[T1]{fontenc}
\usepackage{lmodern}
\usepackage[margin=1in]{geometry}
\usepackage{microtype}
\usepackage{booktabs}
\usepackage{longtable}
\usepackage{array}
\usepackage{xcolor}
\usepackage{fancyhdr}
\usepackage[numbers,sort&compress]{natbib}
\usepackage[hyphens]{url}
\usepackage[hidelinks,breaklinks=true]{hyperref}

% Long DOIs and standards-body URLs are unbreakable tokens; give TeX room to
% resolve them rather than letting them run into the margin.
\setlength{\emergencystretch}{3em}
\makeatletter
\providecommand{\UrlBreaks}{}
\g@addto@macro\UrlBreaks{\do\/\do\-\do\.\do\_\do\=\do\&\do\?}
\makeatother

\definecolor{flagcolor}{RGB}{140,20,20}
\newcommand{\flag}[1]{\textcolor{flagcolor}{\textbf{#1}}}

% Evidence-strength markers used throughout.
\newcommand{\verified}{\textsc{[verified]}}
\newcommand{\secondary}{\textsc{[secondary]}}
\newcommand{\unresolved}{\flag{\textsc{[unresolved]}}}

% A recurring three-part block: the concern, the evidence, the open questions.
\newcommand{\concern}[1]{\par\medskip\noindent\textbf{Concern.} #1}
\newcommand{\evidence}[1]{\par\smallskip\noindent\textbf{What the literature says.} #1}
\newcommand{\investigate}{\par\smallskip\noindent\textbf{Worth investigating.}}

\newcolumntype{L}[1]{>{\raggedright\arraybackslash}p{#1}}

\pagestyle{fancy}
\fancyhf{}
\rhead{\footnotesize Preliminary notes --- not a standard}
\lhead{\footnotesize FL for biological countermeasures}
\cfoot{\thepage}

\title{\textbf{Infrastructure, Data Standards and Data Management for\\
Federated Learning in Biological Countermeasures}\\[0.4em]
\large Preliminary Notes: Concerns and an Investigation Agenda}

\author{}
\date{Evidence current as of 25 September 2026\\[0.3em]
\small Compiled 30 October 2026}

\begin{document}
\maketitle
\thispagestyle{fancy}

\begin{abstract}
\noindent
These are preliminary notes, not a standard. They set out what appears
worth \emph{worrying about} when building federated learning (FL) pipelines
for medical and biological countermeasure work, what the published
literature does and does not establish about each worry, and which
questions seem to need real investigation before anyone commits to an
architecture.

Nothing here should be read as a requirement or a compliance baseline.
Several sections record open questions whose answers would change the
design substantially, and at least one --- whether purely computational
biological design work falls within the scope of current US high-risk
life-sciences research policy (Section~\ref{sec:biosecurity}) --- could
determine whether a programme of this kind is in regulatory scope at all.
The architecture sketch in Section~\ref{sec:architecture} is offered as a
straw man to argue with rather than a recommendation.

Claims carry evidence markers (Section~\ref{sec:howtoread}). The appendix
records what could and could not be verified, and why that limitation is
itself material.
\end{abstract}

\tableofcontents

% ===========================================================================
\section{Purpose and status}

The question these notes respond to is: \emph{how should one think about
infrastructure, data standards and data management so that federated
pipelines for biological countermeasures are safe, secure and reliable?}

They are deliberately framed as concerns rather than answers, for three
reasons.

\begin{itemize}
  \item \textbf{The evidence base is thinner than the enthusiasm.} A 2024
  review of healthcare FL concluded that ``the vast majority are not
  appropriate for clinical use'' \citep{li2024challenges}. Prescribing a
  standard on top of that literature would overstate what is known.

  \item \textbf{Several key questions are genuinely open.} Where the
  literature is contested --- notably adversarial robustness
  (Section~\ref{sec:adversarial}) --- the useful output is a statement of
  the disagreement and what would resolve it, not a picked side.

  \item \textbf{The regulatory ground is moving fast.} Between early 2025
  and mid 2026, the governing US dual-use research policy was rescinded and
  replaced, an AI executive order was rescinded, EU AI Act high-risk
  obligations were deferred, and an export-control rule on AI model weights
  was withdrawn. Section~\ref{sec:regulatory} has a short half-life by
  construction.
\end{itemize}

\paragraph{What these notes are not.} Not a specification, not a controls
baseline, not legal or regulatory advice, and not a claim that federated
learning is the right architecture for any particular programme ---
Section~\ref{sec:alternatives} argues the opposite should be tested first.

\paragraph{Scope boundary worth stating early.} The subject is a pipeline
that produces evidence for human decision-makers. It does not design
biological agents or interventions, does not drive laboratory automation,
and does not make clinical or regulatory determinations. That boundary is
easy to state and harder to establish formally, which is the substance of
Section~\ref{sec:biosecurity}.

% ===========================================================================
\section{How to read this document}
\label{sec:howtoread}

Every substantive claim carries one of three markers reflecting how well it
could actually be checked:

\begin{description}
  \item[\verified] Primary source read first-hand, or DOI resolution
  confirmed directly.
  \item[\secondary] Confirmed via the official page's indexed content plus
  at least two independent corroborating sources, but the primary document
  was not read in full.
  \item[\unresolved] Sources conflicted, or the claim could not be
  corroborated. These are research tasks, not facts, and are collected in
  Section~\ref{sec:agenda}.
\end{description}

The proportion of \secondary{} markers is high, and that is a limitation of
the review rather than of the sources. The appendix explains why.

% ===========================================================================
\section{Starting point: what federation appears to buy}

\subsection{The encouraging evidence}

Biomedical data is rarely poolable: it sits in hospital systems, national
laboratories, public-health agencies, contract research organisations and
commercial sponsors, under incompatible legal bases. Federated learning
moves the model to the data instead \citep{rieke2020future,
kairouz2021advances, sheller2020federated}. Three results suggest this can
produce genuine scientific value rather than only a governance convenience:

\begin{itemize}
  \item \textbf{EXAM} \citep{dayan2021federated} --- 20 hospitals across
  five continents predicting COVID-19 oxygen requirements, reporting average
  AUC above 0.92 and roughly a 38\% average increase in generalisability
  over single-site models. \secondary

  \item \textbf{FeTS} \citep{pati2022federated, pati2023correction} --- the
  largest such study to date at 71 sites across 6 continents, $n = 6{,}314$
  glioblastoma cases, reporting 33\% improvement in delineating the
  surgically targetable tumour and 23\% for complete tumour extent against a
  publicly-trained model. \secondary

  \item \textbf{MELLODDY} \citep{heyndrickx2024melloddy} --- ten competing
  pharmaceutical companies, 2.6 billion confidential activity data points,
  each reporting model improvements without exposing proprietary data.
  Evidence that federation can cross \emph{competitive}, not merely
  institutional, trust boundaries. \secondary
\end{itemize}

\subsection{The counterweight, which deserves equal billing}

These are the published successes, and they required extraordinary
institutional effort. The median outcome looks different:
\citet{li2024challenges} reviewed healthcare FL through May 2024 and
concluded the vast majority of work is not appropriate for clinical use,
citing privacy, generalisation and communication deficiencies. \verified

A frequently repeated claim holds that only about 5.2\% of published
healthcare FL studies represent real clinical deployment. The source for
this could not be reached and the figure should not be used until checked.
\unresolved

\subsection{What inverts relative to centralised machine learning}

\begin{center}
\begin{tabular}{L{0.30\textwidth} L{0.32\textwidth} L{0.26\textwidth}}
\toprule
\textbf{Centralised assumption} & \textbf{Federated reality} &
\textbf{Discussed in} \\
\midrule
One schema, one curation team & $N$ schemas, $N$ curation cultures &
Section~\ref{sec:semantics} \\
Debug by inspecting the data & You never inspect it; you debug through
statistics and contracts & Section~\ref{sec:quality} \\
One trust boundary & Every participant is a dependency \emph{and} a
possible adversary & Section~\ref{sec:adversarial} \\
Availability is one platform's problem & Availability is the product of all
participants' & Section~\ref{sec:reliability} \\
Model output is the only artifact & Gradients, metrics and loss curves are
disclosure channels & Section~\ref{sec:privacy} \\
\bottomrule
\end{tabular}
\end{center}

Federation reduces raw-data movement risk. It does not obviously reduce ---
and may concentrate --- model-mediated disclosure risk, data-quality risk
and dual-use risk.

\subsection{Federation may not be the right answer}
\label{sec:alternatives}

\concern{Choosing FL early, on governance grounds, before testing whether a
simpler architecture would do.}

\evidence{The \textbf{Five Safes} model --- safe projects, people, settings,
data and outputs --- underpins Trusted Research Environments and is the
mature alternative for controlled analysis of sensitive data
\citep{ritchie_fivesafes}. \secondary{} It matters here for a second reason
as well: ``safe outputs'' is \emph{statistical disclosure control}, a
discipline with decades of methodology behind it. Any egress mechanism an FL
platform builds (Section~\ref{sec:privacy}) is a re-invention of that
discipline unless it binds to it deliberately. The EU's European Health Data
Space \citep{ehds2025} is converging on the same shape: computation inside a
certified secure processing environment, with only anonymised or
pseudonymised data leaving.}

\investigate
\begin{itemize}
  \item For the specific scientific question at hand, would a TRE with
  pooled governed access outperform federation on time-to-result, given the
  legal bases actually available?
  \item Is the binding constraint legal, political, or commercial? These
  have different architectural answers, and conflating them tends to
  produce FL by default.
  \item What does the programme lose by never being able to inspect the
  data, and is that cost acknowledged in the project plan?
\end{itemize}

% ===========================================================================
\section{Concerns}

\subsection{Scientific validity and site confounding}
\label{sec:validity}

\concern{A model that performs well across participating sites may be
learning the sites rather than the biology, and conventional validation will
not reveal it.}

\evidence{This is demonstrated, not hypothetical. \citet{zech2018variable}
trained pneumonia classifiers on 158,323 radiographs across three
institutions and found significantly worse performance on outside-hospital
images in 3 of 5 comparisons, because the networks identified hospital
system and department with high accuracy and calibrated to local disease
prevalence. \secondary{} Separately, leakage is a widespread rather than
occasional failure: \citet{kapoor2023leakage} document it across at least
294 papers in 17 disciplines, offering a taxonomy of eight leakage types and
a ``model info sheet'' as a mitigation. \verified{} On the reporting side,
disaggregated evaluation --- performance broken down by demographic,
phenotypic and intersectional groups --- is the defining contribution of
model cards \citep{mitchell2019model} and is routinely omitted. \verified}

\investigate
\begin{itemize}
  \item Can site, instrument, reagent lot and collection period be
  \emph{named} as variables for every participating site? If not, claims
  that the model learned biology rather than logistics are not testable.
  \item How many genuinely held-out sites are available, and who controls
  them? Evaluation that reports to the training team is not evaluation.
  \item Is train/test independence enforced at the site and patient level
  rather than the sample level? Regulatory guidance already expects this
  \citep{gmlp2021}, and it is the split most often got wrong.
  \item What is the pre-registered evaluation plan, and what would count as
  a negative result?
\end{itemize}

\subsection{Data semantics: what a row means at every site}
\label{sec:semantics}

\concern{The hardest part of federated work is not the learning. It is
agreeing what a field means across institutions that never compare raw
records --- and the failure mode is silent.}

\evidence{A substantial standards ecosystem exists and should be bound to
rather than reinvented. Table~\ref{tab:standards} records what could be
established about each, including version details that matter in practice.}

\begin{center}
\small
\begin{longtable}{L{0.20\textwidth} L{0.22\textwidth} L{0.46\textwidth}}
\caption{Standards worth binding to, with version notes. Verified where the
domain was reachable; most standards-body sites were blocked from the review
environment (see appendix).}
\label{tab:standards}\\
\toprule
\textbf{Domain} & \textbf{Standard} & \textbf{Notes} \\
\midrule
\endfirsthead
\toprule
\textbf{Domain} & \textbf{Standard} & \textbf{Notes} \\
\midrule
\endhead
\bottomrule
\endfoot
Clinical / EHR exchange & HL7 FHIR \citep{hl7fhir} &
R5 (v5.0.0, 2023) is current published; R6 is in normative ballot only. But
R4 (v4.0.1) remains the deployment reality for US Core and EHDS work ---
choose deliberately. \secondary \\
\addlinespace
Observational analytics & OMOP CDM \citep{omopcdm} &
v5.5 current, a non-breaking superset of v5.4. v6.0 exists but is
\emph{unsupported} by the OHDSI tool stack; never cite it as latest.
\verified \\
\addlinespace
Clinical trial data & CDISC SDTM, ADaM \citep{cdisc} &
SDTM is as-collected tabulation, ADaM analysis-ready with traceability.
Version numbers \unresolved. \\
\addlinespace
Nonclinical toxicology & CDISC SEND \citep{cdisc} &
SEND is SDTM \emph{applied to the nonclinical domain} --- animal toxicology,
not clinical. Include only if preclinical safety is in scope. \secondary \\
\addlinespace
Lab observations & LOINC \citep{loinc} &
Names what was measured, not the result. Free licence. \secondary \\
\addlinespace
Clinical concepts & SNOMED CT \citep{snomedct} &
\flag{Requires a licence} --- free in member countries via national release
centres, paid otherwise. \secondary \\
\addlinespace
Units & UCUM \citep{ucum} &
``Unified \emph{Code} for Units of Measure'' --- singular. \secondary \\
\addlinespace
Phenotype exchange & GA4GH Phenopackets v2 \citep{phenopackets} &
Published as ISO 4454:2022, the first GA4GH standard published by ISO.
\secondary \\
\addlinespace
Data resolution & GA4GH DRS v1.5.0 \citep{ga4ghdrs} &
An ID-to-access-URL resolution layer, \emph{not} a transfer protocol and not
a search API. \verified \\
\addlinespace
Researcher authorisation & GA4GH AAI v1.2 \citep{ga4ghaai} and Passport
v1.2 \citep{ga4ghpassport} &
\flag{Two distinct specifications} from two work streams. AAI is the OIDC
protocol; Passport is the token format. Writing ``Passports/AAI'' as one
thing is imprecise. \verified \\
\addlinespace
Imaging de-identification & DICOM PS3.15 Annex E \citep{dicomps315} &
E.2 Basic Profile plus named Options; Table E.1-1 is the per-attribute
action table. \flag{The standard itself states the profiles do not guarantee
de-identification and do not replace a de-identification process.}
Burned-in pixel data, private tags and free-text descriptors are the
residual-risk vectors. \secondary \\
\addlinespace
Experiment metadata & ISA framework \citep{isatools} &
Domain-neutral community framework, ELIXIR Recommended Interoperability
Resource. Not a ratified SDO standard --- describe it accurately.
\secondary \\
\addlinespace
Standards discovery & FAIRsharing \citep{mcquilton2016biosharing} &
The live successor to MIBBI \citep{taylor2008mibbi} via BioSharing. MIBBI no
longer exists as an active registry. \secondary \\
\addlinespace
Organism identity & NCBI Taxonomy \citep{schoch2020ncbi} &
The nomenclature repository for INSDC. \emph{Not} an authoritative
phylogenetic or nomenclatural authority --- NCBI says so itself.
\secondary \\
\addlinespace
Provenance & W3C PROV-O \citep{w3cprov} &
Recommendation since 2013, stable. Cite PROV-O specifically for the RDF
ontology. \secondary \\
\addlinespace
Catalogue metadata & W3C DCAT 3 \citep{w3cdcat3} &
Recommendation, August 2024, built on PROV-O terms. The health profile
HealthDCAT-AP is separate and evolving --- do not conflate. \secondary \\
\addlinespace
Findability principles & FAIR \citep{wilkinson2016fair} &
See the caveat below: FAIR is not a standard. \verified \\
\end{longtable}
\end{center}

\paragraph{Three corrections worth recording, because all three are easy
mistakes.}

\begin{itemize}
  \item \flag{MIAPPE does not belong in a human biomedical standards set.}
  It is the \emph{Minimum Information About a Plant Phenotyping Experiment}
  \citep{miappe}; its entities are germplasm, growth facility, field plot
  layout and agronomic treatment. An earlier draft of these notes cited it.
  That was wrong, and it is the kind of error that signals a standards list
  assembled without domain checking. \secondary

  \item \flag{MINSEQE is orphaned.} The FGED Society dissolved in September
  2021. It is citable as a historical checklist only, never as a live
  standard with a governing body. \secondary

  \item \flag{``FAIR compliance'' is not a meaningful claim.} FAIR
  \citep{wilkinson2016fair} has no conformance test, no version and no
  certifying body. The defensible phrasing is that a platform
  \emph{implements} the FAIR principles via DCAT, PROV and persistent
  identifiers. Note also that \emph{Accessible} does not mean \emph{open}
  --- FAIR explicitly accommodates authenticated, controlled access, which
  is precisely why GA4GH Passports and AAI are its right companions for
  human data.
\end{itemize}

\paragraph{The mechanism that seems to matter most.} A versioned,
machine-readable \emph{feature contract} per model family --- authored
centrally, enforced locally --- specifying for each field: logical name,
semantic type and unit; source code system with version; permissible values
and the action on violation; whether missingness is informative and how it
is encoded; derivation as executable code rather than prose; and
sensitivity class. Harmonisation then runs at the edge, so mapping logic
stays with the people who understand the source system, and mapping bugs
surface in conformance reports rather than requiring access to raw records.

\investigate
\begin{itemize}
  \item Which standard versions will actually be bound, and is every coded
  field carrying \texttt{(system, version, code)}?
  \item What happens to unmapped values? Silent coercion is the failure
  mode; an explicit unmapped queue with the original string retained
  locally is the alternative.
  \item How are contract versions enforced across a round --- can a round
  silently mix versions?
  \item Who owns each site's mapping specification, and is an unreviewed
  mapping change treated as the production change it is?
  \item Does SNOMED CT licensing actually cover every participating
  jurisdiction? This is a real blocker that tends to surface late.
\end{itemize}

\subsection{Data quality without being able to look}
\label{sec:quality}

\concern{Quality must be asserted, tested remotely and attested --- and
there is a strong temptation to invent a bespoke taxonomy for this.}

\evidence{A canonical framework exists. \citet{kahn2016harmonized} is the
harmonised terminology for secondary use of EHR data, and it has
\emph{two} axes, the second of which is routinely dropped:

\begin{description}
  \item[Categories] \textbf{Conformance} (value, relational,
  computational); \textbf{Completeness}; \textbf{Plausibility} (uniqueness,
  atemporal, temporal). Note the asymmetry --- conformance and plausibility
  have subcategories, completeness does not.
  \item[Contexts] \textbf{Verification} against internal expectations and
  metadata; \textbf{Validation} against an external gold standard or
  real-world referent.
\end{description}

\verified{} The OHDSI Data Quality Dashboard
\citep{blacketer2021increasing, ohdsi_dqd} operationalises exactly this
taxonomy against OMOP as 24 parameterised check types instantiated as more
than 3,000 individual checks, each tagged verification or validation. (Those
two numbers describe templates and instantiations respectively and are often
conflated.) \verified}

\paragraph{But federated quality assessment has no agreed standard.} Kahn
supplies terminology, the DQD supplies an OMOP implementation, and the
operational patterns come from FDA Sentinel (a coordinating centre
distributes check programs, sites execute locally, results return) and
PCORnet (which adapted Sentinel's approach to assess ``foundational data
quality'' --- fitness across a research portfolio rather than per study).
The standard critique is that these are \emph{top-down}: sites must satisfy
centrally defined targets rather than assess in local context, and
cross-network comparability remains open. \secondary

\investigate
\begin{itemize}
  \item Can the programme adopt Kahn's terminology directly, including the
  verification/validation axis, rather than growing a parallel scheme?
  \item What is the actual gate policy --- are failing sites excluded with a
  stated reason, or quietly down-weighted? The latter hides the problem in
  the model.
  \item How is cross-site duplication \emph{risk} estimated without
  exchanging identifiers?
  \item How is label quality established: provenance, adjudication
  procedure, inter-rater agreement?
  \item Is there an explicit leakage check for fields encoding the outcome
  or post-outcome information \citep{kapoor2023leakage}?
  \item On deletion: does it propagate to snapshots, caches, checkpoints and
  derived artifacts, with a per-participant attestation? And does consent
  withdrawal imply retraining? Removing one subject's influence from a
  trained model retroactively is expensive and sometimes infeasible, so the
  answer belongs in the data-use agreement rather than in an incident.
\end{itemize}

\subsection{Privacy: stating precisely what each mechanism guarantees}
\label{sec:privacy}

\concern{Privacy mechanisms are often described in ways that overstate what
they deliver, and federation itself is frequently assumed to remove a
programme from data-protection scope.}

\evidence{Three findings deserve to be stated exactly.

\textbf{Secure aggregation} \citep{bonawitz2017practical} gives a
communication-efficient, dropout-robust protocol by which a server computes
a \emph{sum} over a cohort without learning any individual contribution.
\secondary{} Its limit: it protects the individual update, not the
aggregate. Anything inferable from the sum remains inferable. Secure
aggregation and differential privacy are complements, not substitutes.

\textbf{Gradient inversion is more feasible than often assumed.}
\citet{geiping2020inverting} reconstruct images at high resolution from
parameter gradients, report that input to a fully connected layer is
analytically reconstructible independent of the rest of the architecture,
and state that averaging gradients over several iterations or several images
does \emph{not} protect privacy. \secondary{} An earlier draft of these
notes offered large cohort sizes as a defence against inversion; that was
not supportable.

\textbf{The unit of privacy is where $\varepsilon$ reporting goes wrong.}
Example-level and client-level differential privacy give materially
different guarantees. In a multi-hospital setting the entity needing
protection is usually the \emph{patient}, who may contribute many records,
so a single-record definition degrades accordingly. An $\varepsilon$
reported without its unit, accounting method and composition assumptions is
not a privacy claim. \secondary

End-to-end integration has been demonstrated: PriMIA
\citep{kaissis2021endtoend} combines differentially private, securely
aggregated federated learning with encrypted inference on paediatric chest
radiographs, reporting performance comparable to locally trained models and
empirical resistance to a gradient-based inversion attack; see also
\citet{kaissis2020secure}. \secondary}

\paragraph{Data-protection scope.} Gradients, embeddings and weights derived
from personal data are frequently still personal data under GDPR
\citep{gdpr, edpb28_2024}; federation does not automatically take a
programme out of scope. Federated consortia also usually create
\emph{joint controllership} rather than processorship, which is an Article 26
arrangement best settled before the first training round. \secondary

\investigate
\begin{itemize}
  \item What exactly is the disclosure-control policy on egress --- artifact
  types, shapes, minimum cell sizes, suppression rules --- and does it bind
  to statistical disclosure control practice \citep{ritchie_fivesafes}
  rather than being bespoke?
  \item What is the declared unit of privacy, and who enforces the budget?
  A ledger the platform enforces behaves differently from a convention.
  \item Are per-site metrics, loss curves and error analyses treated as
  disclosure channels, or do they circulate informally? The latter is
  common and rarely reviewed.
  \item Would a dedicated review of the DP accounting literature change the
  parameter choices? These notes did not examine it in depth and should not
  be relied on for it. \unresolved
  \item Confidential computing is attractive but not a clean answer:
  side-channel attacks are the primary documented threat to TEEs and enclave
  memory limits constrain ML workloads \citep{teesok2022}. Is it being used
  as defence in depth or as the privacy argument? \secondary
\end{itemize}

\subsection{Adversarial robustness, where the literature genuinely disagrees}
\label{sec:adversarial}

\concern{The obvious reading of the attack literature is that federated
learning is dangerously fragile. The obvious reading of the critique
literature is that it is fine. Both readings are available, and an earlier
draft of these notes picked a side without acknowledging the dispute.}

\evidence{Laid out in sequence:

\begin{itemize}
  \item \textbf{Attacks succeed under their stated assumptions.}
  \citet{bagdasaryan2020backdoor} show model-replacement backdoors reaching
  100\% backdoor accuracy in a single round, with an attacker controlling
  fewer than 1\% of participants able to prevent the model from unlearning
  the backdoor. \secondary

  \item \textbf{Defences have formal but bounded guarantees.} Krum
  \citep{blanchard2017machine} and coordinate-wise median / trimmed mean
  \citep{yin2018byzantine}, the latter with order-optimal statistical error
  rates under subexponential gradient noise, are the canonical robust
  estimators. \secondary

  \item \textbf{The threat models were largely unrealistic.}
  \citet{shejwalkar2022back} systematise the space and conclude that prior
  work assumed impractically high compromised-client fractions and adversary
  capability, finding that ``FL is highly robust in practice even when using
  simple, low-cost defenses.'' \verified

  \item \textbf{Defence evaluation remains weak.}
  \citet{sok2025benchmarking} observe that defences are typically evaluated
  in isolation against limited attack strategies, and benchmark 15 attacks
  against 17 defences across 2,040 settings. \secondary
\end{itemize}}

\paragraph{Why the disagreement may not resolve in this setting.}
\citet{shejwalkar2022back} ground their reassurance in \emph{cross-device}
production FL with large participant populations. A biomedical consortium is
\emph{cross-silo}: few participants, each holding large influence, where a
single compromised site is a far larger fraction of the update than in the
setting that study examined. The reassurance may simply not transfer. What
does appear to transfer, because it does not depend on out-detecting a
clever adversary, is contractual and identity vetting at enrollment, bounds
enforcement on update norms and per-cohort contribution, per-site influence
monitoring, and independent held-out evaluation.

\investigate
\begin{itemize}
  \item Does any published result characterise poisoning robustness in the
  cross-silo, few-participant regime specifically? If not, this is a
  research gap the programme may have to fill itself rather than assume
  away.
  \item What bound on any single participant's per-round influence is
  acceptable, and is it enforced by construction or detected after the
  fact?
  \item Is there an isolation procedure for a misbehaving participant, and
  has it been \emph{rehearsed}? An unrehearsed kill switch is a hypothesis.
  \item For threat-model vocabulary, NIST AI 100-2e2025
  \citep{nistai1002} organises attacks by lifecycle stage and attacker
  goal, capability and knowledge, with an attack/mitigation index. Is the
  programme's threat model written in a shared vocabulary or an ad hoc one?
  \secondary
  \item Supply chain: SLSA \citep{slsa} asserts build-platform integrity and
  in-toto \citep{intoto} supplies the attestation envelope. The interesting
  question is whether in-toto predicates can be extended beyond build
  provenance to \emph{dataset}, \emph{model-weight} and
  \emph{evaluation-result} attestation, so a model carries a verifiable
  chain back to its training corpus and its review.
\end{itemize}

\subsection{Reliability when partial participation is normal}
\label{sec:reliability}

\concern{Sites will be unavailable, slow, mid-migration or failing quality
gates. A design that treats full participation as the default will stall;
but the dominant risk may not be availability at all.}

\evidence{Statistical heterogeneity, not downtime, appears to be the
dominant performance risk: non-IID label and feature distributions across
sites cause client drift and unstable convergence, and quantity imbalance
compounds it. \secondary{} The practical implication is that the engineering
needed for partial participation --- quorum rounds, bounded staleness,
per-round checkpointing, idempotent retryable stages, cohort-weighted
aggregation with caps --- is necessary but will not by itself produce a
model that transfers.}

\investigate
\begin{itemize}
  \item What participation floor constitutes a valid round, and is who
  participated recorded per round?
  \item What is the straggler and deadline policy, and is it in
  configuration or decided case by case?
  \item How much work does a coordinator failure cost, and has recovery been
  tested by drill rather than argued on paper?
  \item Which service levels should have \emph{no} error budget? Integrity
  and privacy objectives are plausible candidates --- they are not
  tradeable against delivery speed in the way latency is.
  \item Are adversarial and chaos drills scheduled --- dropping a site
  mid-round, corrupting a delta, replaying a stale update, simulating
  coordinator loss?
\end{itemize}

% ===========================================================================
\section{Regulatory and governance currency}
\label{sec:regulatory}

\concern{Compliance artifacts built on remembered rather than checked policy
will be wrong, because an unusual amount changed in 2025--2026. This section
is the most perishable in the document.}

\subsection{Medical device and pharmaceutical regimes}

\begin{itemize}
  \item \textbf{The FDA PCCP guidance is final} \citep{fdapccp2024} and is
  arguably the single most consequential document for pipeline
  architecture. A Predetermined Change Control Plan authorises pre-specified
  modifications without a new marketing submission, requiring a bounded
  Description of Modifications, a Modification Protocol (data management,
  re-training, performance evaluation and acceptance criteria) and an Impact
  Assessment. The design consequence: a retraining loop, drift monitoring,
  dataset versioning and acceptance thresholds must be \emph{writable into a
  Modification Protocol up front}, or each model refresh becomes a new
  submission. \secondary

  \item \textbf{The FDA AI device lifecycle guidance remains draft}
  \citep{fdalifecycle2025} --- issued January 2025, not finalised. Not
  binding. \secondary

  \item \textbf{The FDA AI-in-drug-development guidance}
  \citep{fdaaidrug2025} introduces a seven-step risk-based Credibility
  Assessment Framework anchored on Context of Use, implying a credibility
  dossier \emph{per context of use} rather than per model. Whether a final
  version issued in 2026 is \unresolved.

  \item \textbf{FDA and EMA published joint AI principles in January 2026}
  \citep{fdaema2026} --- ten principles, the first such alignment for
  medicines. \secondary

  \item \textbf{GMLP principles} \citep{gmlp2021} are what reviewers
  actually apply. Principle 4 (independent training and test datasets)
  dictates splitting at site and patient level; principle 7 evaluates the
  human--AI team rather than the model alone. \secondary

  \item \textbf{Lifecycle and quality standards are stable}, which is worth
  knowing precisely because other things are not: IEC 62304 Edition 1.1
  remains current with Edition 2 unpublished \citep{iec62304}; ISO 14971
  \citep{iso14971} was reconfirmed in 2025 and its post-production feedback
  loop is where drift monitoring legally attaches; ISO 13485
  \citep{iso13485} was reconfirmed in late 2025 and is incorporated by
  reference into the FDA QMSR effective February 2026. IEC 62304's
  \emph{SOUP} provisions are where an ML framework, CUDA stack and any
  pretrained weights land --- a foundation model in a Class C device is SOUP
  with an awkward anomaly history. \secondary

  \item \textbf{ALCOA+ maps onto technical controls} more cleanly than its
  reputation suggests \citep{mhra2018dataintegrity, pics041}:
  attributable/contemporaneous to a signed, server-timestamped immutable
  audit log; original to retaining raw source plus a verified true copy
  rather than only derived features; enduring/available to readability over
  a retention period that for biologics may outlive the ML framework
  itself. \secondary

  \item \flag{Draft EU GMP Annex 22 on artificial intelligence}
  \citep{gmpannex22draft} went to consultation in July 2025, closed in
  October 2025, with final text targeted for Q4 2026 --- so not in force. As
  drafted it would restrict AI in GMP-critical applications to
  \emph{static, deterministic models}, excluding dynamic and
  continuously-learning models. A manufacturing decision path that assumes a
  continuously-learning model is therefore at architectural risk.
  \secondary
\end{itemize}

\subsection{EU data and AI law}

\begin{itemize}
  \item \flag{EU AI Act high-risk obligations were deferred.} Regulation
  (EU) 2024/1689 \citep{euaiact2024} was amended by Regulation (EU)
  2026/1744 \citep{euomnibus2026}, in force 27 July 2026 --- six days before
  the original deadline. Annex III standalone high-risk moved from 2 August
  2026 to \textbf{2 December 2027}; Annex I, which covers AI in MDR/IVDR
  devices, moved to \textbf{2 August 2028}. Prohibited practices,
  general-purpose AI obligations and Article 50 transparency duties were
  \emph{not} deferred. \verified{} Any analysis written before July 2026 that
  assumes an August 2026 bite is now wrong.

  \item Most clinical AI reaches high-risk via \emph{Annex I rather than
  Annex III}, because an AI-enabled medical device already requires
  third-party conformity assessment. Articles 9 and 17 map substantially
  onto ISO 14971 and ISO 13485; \textbf{Article 10 on data and data
  governance has no clean MDR analogue and is the genuine new work};
  Article 12 automatic logging is a hard architectural requirement.
  \secondary

  \item \flag{The GDPR ``Data Omnibus'' was not adopted.} Only the AI
  Omnibus passed. Proposed GDPR amendments --- including making research
  processing always purpose-compatible --- remain in negotiation, and the
  EDPB/EDPS joint opinion of February 2026 is strongly critical. Designing
  to the proposed text would be a mistake. \secondary

  \item The \textbf{European Health Data Space} \citep{ehds2025} is in force
  since March 2025, with general application March 2027 and the
  secondary-use regime from \textbf{March 2029}: access via national Health
  Data Access Bodies issuing permits, inside a secure processing
  environment, with only anonymised or pseudonymised data leaving. Building
  for that access pattern now is not premature. \secondary
\end{itemize}

\subsection{Security and AI-risk frameworks}

\begin{itemize}
  \item \textbf{NIST SP 800-53} \citep{nist80053} --- pin the
  \emph{release}, not only the revision. Release 5.2.0 (August 2025) adds
  controls on software resiliency by design, developer testing, secure
  update and patch deployment, and software integrity. \secondary
  \item \textbf{NIST CSF 2.0} \citep{nistcsf2} adds the GOVERN function and
  strengthened supply-chain outcomes, which is what auditors use to ask who
  owns model-risk decisions. \secondary
  \item \textbf{Zero trust} \citep{nist800207} --- treat each training job,
  feature store, model registry and inference endpoint as a separately
  authenticated subject rather than relying on network perimeter.
  \secondary
  \item \flag{The NIST AI RMF is under active revision}
  \citep{nistairmf} --- the NIST AI Resource Center carries a notice to that
  effect, with no 1.1 or 2.0 published. \verified{} Pin the version and date
  in any compliance artifact. The Generative AI Profile
  \citep{nistai6001} identifies CBRN information and capability among twelve
  risk areas. \secondary
  \item \textbf{ISO/IEC 42001} \citep{iso42001} is a \emph{management
  system} standard, not a technical AI-safety standard. It provides a
  governance wrapper; it will not satisfy EU AI Act technical requirements,
  FDA credibility expectations or any biosecurity control on its own.
  \flag{ISO/IEC 27701:2025} \citep{iso27701} is now standalone and
  separately certifiable, with scope expanded to AI-related processing and
  health data. \secondary
\end{itemize}

% ===========================================================================
\section{Biosecurity and dual-use: the largest open question}
\label{sec:biosecurity}

\concern{The governing US policy changed fundamentally in 2026, and whether
a computational pipeline of this kind is in scope at all is unresolved.}

\evidence{The 2024 DURC/PEPP policy \citep{durcpepp2024} was overtaken by
Executive Order 14292 \citep{eo14292}, signed 5 May 2025 --- one day before
that policy's effective date --- which paused federal funding for dangerous
gain-of-function research and directed its revision or replacement.

The replacement is the \textbf{US Government Policy for Stopping High-Risk
Life Sciences Research} \citep{usgstopping2026}, approved 20 July 2026 and
released 28 July 2026. \verified{} It replaces the 2024 DURC/PEPP policy and
the earlier DURC and P3CO policies; shifts from a list-based to a risk-based
framework and from ``oversee and mitigate'' to \emph{prohibit and stop},
establishing federal funding prohibitions rather than only review
requirements; introduces two categories, \textbf{Dangerous
Gain-of-Function (DGOF)} research and \textbf{International Research of
Concern (IROC)}; restricts federally supported life-sciences research
conducted \emph{outside} the United States; and imposes expanded review,
certification, monitoring, reporting and enforcement obligations on all
institutions receiving federal life-sciences funding.}

\paragraph{The question that matters most, and is not answered here.}
Whether purely computational or in-silico biological design work falls
within ``life sciences research'' under this policy \emph{determines whether
a pipeline of this kind is in regulatory scope at all}. The policy text
could not be retrieved from the review environment, so the DGOF and IROC
definitions are \unresolved. This should be the first thing read in full.
Institutions are meanwhile acting beyond the federal floor --- the
University of California mandated a pause on work meeting the DGOF
definition regardless of funding source, effective May 2026. \secondary

\paragraph{Other changes that overturn common assumptions.}

\begin{itemize}
  \item \flag{EO 14110 was rescinded} \citep{eo14110} in January 2025;
  EO 14179 and the July 2025 AI Action Plan \citep{aiactionplan2025} are
  operative. Importantly, \emph{artifacts produced under} EO 14110 ---
  including the NIST AI RMF and agency AI policies --- survive its
  rescission. \secondary

  \item \flag{Nucleic acid synthesis screening is in limbo.} The April 2024
  OSTP framework \citep{nucleicacid2024} made provider attestation a
  condition of federal research funding, but EO 14292 directed its revision
  or replacement within 90 days, a deadline that passed in August 2025 with
  no replacement published. The Sequences of Concern definition is expected
  to widen. If any part of a pipeline emits sequences for synthesis, the
  implication is to build provider-attestation verification and an
  order-level screening audit trail now, and to treat the SOC definition as
  a configuration value rather than a design assumption. \secondary

  \item \textbf{AI biological-capability evaluation is by agreement, not
  regulation.} There is no US statutory requirement to evaluate a model for
  biological uplift before release; the binding obligations sit in the EU AI
  Act general-purpose AI systemic-risk regime and in provider-internal
  frameworks. If a pipeline consumes a frontier model, the enforceable
  safeguards sit with the provider under EU law and with the consumer
  contractually. \secondary

  \item \textbf{Where to put the hardest controls.} The NASEM consensus
  study \citep{nasem2025ai} assesses AI-enabled biological design tools and
  finds that AI is also a biosecurity \emph{strengthener}. Its most
  architecturally useful conclusion is that risk concentrates at the
  AI-to-wet-lab interface --- synthesis and physical execution --- rather
  than in the model itself, which argues for concentrating controls at the
  design-to-order boundary rather than attempting to gate the model.
  \secondary

  \item \flag{The NIH Guidelines are being replaced.} A draft NIH Biosafety
  Policy for Research Involving Biohazards \citep{nihdraftbiosafety2026} was
  released in August 2026 and would replace the NIH Guidelines
  \citep{nihguidelines} entirely, widening scope from recombinant and
  synthetic nucleic acids to biohazards generally with tiered risk-based
  oversight. The comment period closed 19 October 2026; the Guidelines
  remain operative until the replacement is final. Separately, IBC minutes
  from meetings on or after 1 June 2025 must be posted publicly --- so
  institutional deliberations about a programme should be assumed public.
  \secondary
\end{itemize}

\paragraph{Export control, briefly.} Biotech equipment controls tightened in
January 2025 with new ECCNs 3A069 and 3E069 \citep{bis3a069}, the rule's
rationale explicitly citing the combination of such equipment with AI and
biological design tools. Conversely, \flag{there is currently no general US
export control on AI model weights} --- the AI Diffusion Rule was rescinded
in May 2025 \citep{bisdiffusionrescind} --- so a risk register assuming
otherwise is wrong. The trap is \emph{deemed export}: releasing controlled
technology, including source code and technical data, to a foreign national
inside the United States is a licensable event, which means repository,
model-weight and sequence-data access control is an export-control surface
and not only an HR process. \secondary

\paragraph{Gate stacking.} A countermeasure programme typically passes four
to five independent review gates with different owners: institutional
biosafety; high-risk/dual-use review; IRB for human data (with single-IRB
requirements for NIH-funded multi-site human-subjects research); select
agent registration where applicable; and export-control review. The
dual-use gate is the one most often discovered late.

% ===========================================================================
\section{Documentation and reporting}

\concern{There is a temptation to assume a federated-learning reporting
standard exists. It does not.}

\evidence{No EQUATOR-registered consensus reporting guideline is specific to
federated learning. \secondary{} The stage-appropriate general guidelines
are summarised in Table~\ref{tab:reporting}; confusing SPIRIT-AI (protocol)
with CONSORT-AI (report) is the most common error with that pair.}

\begin{table}[htbp]
\centering
\caption{Stage-appropriate reporting guidelines. No equivalent exists that is
specific to federated learning.}
\label{tab:reporting}
\begin{tabular}{L{0.34\textwidth} L{0.56\textwidth}}
\toprule
\textbf{Lifecycle stage} & \textbf{Guideline} \\
\midrule
Model development / validation & TRIPOD+AI \citep{collins2024tripodai} ---
unifies regression and ML, adds fairness and open-science items \\
Early-stage live clinical evaluation & DECIDE-AI \citep{vasey2022decideai}
--- human factors and clinician--AI interaction, not accuracy \\
Trial protocol & SPIRIT-AI \citep{cruzrivera2020spiritai} \\
Trial report & CONSORT-AI \citep{liu2020consortai} \\
Medical imaging AI & CLAIM 2024 Update \citep{tejani2024claim}; cite this,
not the 2020 original \citep{mongan2020claim} \\
Baseline minimum information & MINIMAR
\citep{hernandezboussard2020minimar} --- a proposal, not a
Delphi-consensus guideline \\
\bottomrule
\end{tabular}
\end{table}

\noindent
All DOIs in Table~\ref{tab:reporting} resolve to the correct publisher
paths. \verified{} Supplement with FUTURE-AI
\citep{lekadir2025futureai} --- 118 experts across 51 countries, a two-year
modified-Delphi process, six principles (fairness, universality,
traceability, usability, robustness, explainability) expanded into 28 best
practices --- noting it is a trustworthiness and deployment framework rather
than a reporting checklist. \verified

\investigate
\begin{itemize}
  \item Regardless of which guideline applies, federated studies arguably
  ought to report site count, per-site sample sizes and distributions,
  aggregation strategy, communication rounds, privacy mechanism with unit of
  privacy and $\varepsilon$ accounting, and \emph{per-site rather than
  pooled} performance. Is there appetite to standardise that set?
  \item Per artifact, three documents seem worth insisting on: a datasheet
  per dataset snapshot \citep{gebru2021datasheets}, a model card with
  disaggregated evaluation per model \citep{mitchell2019model}, and a model
  info sheet for leakage \citep{kapoor2023leakage}. Datasheets are
  qualitative by design with no conformance criteria, so pairing them with
  machine-readable DCAT/PROV records is probably necessary.
\end{itemize}

% ===========================================================================
\section{An architecture sketch, offered as a straw man}
\label{sec:architecture}

The following is a plausible shape, included so that it can be argued with.
It is not a recommendation, and the concerns above are more important than
this diagram.

\begin{verbatim}
  +---------------------------------------------------------------------+
  | GOVERNANCE PLANE                                                    |
  |  enrollment & identity . data-use agreements . approvals            |
  |  privacy-budget ledger . model registry . immutable audit log       |
  +---------------------------------------------------------------------+
                  | policy decisions, attestations, approvals
  +---------------------------------------------------------------------+
  | ORCHESTRATION PLANE   (coordinator - untrusted with raw data)       |
  |  round scheduling . participant selection . secure aggregation      |
  |  bounds enforcement . robust aggregation . checkpointing            |
  +---------------------------------------------------------------------+
     ^ signed deltas / encrypted shares        v signed global model
  +---------------------------------------------------------------------+
  | PARTICIPANT PLANE   (per site, inside the site's boundary)          |
  |  ingest -> harmonization -> QC gates -> snapshot -> local training  |
  |  local eval . egress disclosure control . local audit log           |
  +---------------------------------------------------------------------+
  +---------------------------------------------------------------------+
  | EVALUATION PLANE   (independent of training)                        |
  |  held-out site benchmarks . drift & subgroup monitors . red team     |
  +---------------------------------------------------------------------+
  +---------------------------------------------------------------------+
  | CONSUMPTION PLANE                                                   |
  |  serving . decision support . reporting . feedback capture          |
  +---------------------------------------------------------------------+
\end{verbatim}

The propositions embedded in it, each of which is contestable:

\begin{enumerate}
  \item The coordinator is untrusted with data and should be designed as
  though it will be breached.
  \item Egress is an enforced chokepoint implementing statistical
  disclosure control, default-deny.
  \item Evaluation does not report to training --- different owners,
  different credentials, separate held-out data.
  \item Reproducibility is infrastructure: content-addressed artifacts,
  pinned digests, a run manifest per round.
  \item Governance is enforced in code paths, since an approval that exists
  only in a document is not a control.
\end{enumerate}

% ===========================================================================
\section{Consolidated investigation agenda}
\label{sec:agenda}

Ordered by how much the answer would change the design.

\begin{enumerate}
  \item \flag{Read the July 2026 \emph{Stopping High-Risk Life Sciences
  Research} policy in full} \citep{usgstopping2026} --- specifically the
  DGOF and IROC definitions, and whether in-silico design is in scope. This
  determines scope exposure more than anything else here, and is currently
  unresolved.

  \item \textbf{Decide federation versus a Trusted Research Environment}
  on evidence rather than by default (Section~\ref{sec:alternatives}),
  including what is lost by never inspecting the data.

  \item \textbf{Establish whether poisoning robustness has been
  characterised in the cross-silo, few-participant regime}
  (Section~\ref{sec:adversarial}). If not, it is a gap to fill rather than
  assume away.

  \item \textbf{If a regulated device is plausible, draft the PCCP
  Modification Protocol early} \citep{fdapccp2024} --- it constrains
  retraining, drift monitoring, dataset versioning and acceptance
  thresholds downstream.

  \item \textbf{Commission a dedicated review of DP accounting and the unit
  of privacy.} These notes establish that the question matters but do not
  settle the parameter choices.

  \item \textbf{Verify CDISC version numbers} directly
  (Table~\ref{tab:standards}) --- currently low confidence.

  \item \textbf{Confirm whether the FDA January 2025 AI-in-drug-development
  draft was finalised} \citep{fdaaidrug2025} --- sources conflicted.

  \item \textbf{Verify SLSA v1.2 level semantics} first-hand
  \citep{slsa}.

  \item \textbf{Track the replacement nucleic acid synthesis screening
  framework} \citep{nucleicacid2024} and keep the Sequences of Concern
  definition configurable.

  \item \textbf{Check the ``5.2\% of healthcare FL studies are real
  deployments'' figure} before it is used anywhere.
\end{enumerate}

\paragraph{Signals that a programme is drifting.} Raw data requested ``just
to debug''; per-site metrics circulating informally; mapping changes without
review; a benchmark the training team can see; an isolation procedure never
exercised; a model card written after the model shipped; privacy parameters
chosen to hit an accuracy target; a dual-use determination deferred to
``later''.

% ===========================================================================
\appendix
\section{Verification limits}

These notes were compiled in an environment whose network egress proxy
blocked most standards-body, regulator and publisher domains --- including
\texttt{nature.com}, PubMed Central, ScienceDirect, \texttt{hl7.org},
\texttt{cdisc.org}, \texttt{dicom.nema.org}, \texttt{w3.org},
\texttt{nist.gov}, \texttt{fda.gov}, \texttt{whitehouse.gov},
\texttt{eur-lex.europa.eu} and \texttt{iso.org}. Reachable:
\texttt{arxiv.org}, \texttt{doi.org} for resolution,
\texttt{raw.githubusercontent.com}, and search indexes.

Consequently most citations are marked \secondary: DOIs were resolution-
tested and official URLs were drawn from official-domain search hits rather
than constructed, but the primary documents were largely not read in full.
Two claims with the highest stakes --- the EU AI Act deferral
\citep{euomnibus2026} and the existence and effect of the 2026 high-risk
life sciences policy \citep{usgstopping2026} --- were verified
independently. The full text of the latter was still not obtainable, which
is why Section~\ref{sec:biosecurity} leaves its central question open rather
than inferring an answer.

Widening the network allowlist would materially raise confidence in
Table~\ref{tab:standards} and Sections~\ref{sec:regulatory}
and~\ref{sec:biosecurity}. Until then, the evidence markers should be taken
literally: \secondary{} means corroborated but not read, and \unresolved{}
means do not rely on it.

% ===========================================================================
\bibliographystyle{plainnat}
\bibliography{references}

\end{document}
